\section{QFT toys}
\label{ch:qft-toys}

This appendix collects three field-theory models in which the Hubbard--Stratonovich step of the construction can be done in a few lines, with every coefficient explicit.
The first block is the paper's appendix on $N$ scalar fields with a quartic interaction, reproduced in full: the auxiliary field carries a pair of indices as $\Sigma_2$ carries two strings, and a degenerate coupling matrix shows how the kernel condition and the multiplier on the null directions arise in finite dimension.
The second block, from the older overview note, gives the strictification of the quartic Yang--Mills vertex in the colour--kinematics-dual form and the open-string field theory with stubs, the closest string-theoretic precedent of an auxiliary field with a pole-free kinetic term, together with the list of what does and does not carry over to closed strings.
The third block, from the companion note on pairings and gradings, redoes the Yang--Mills elimination step by step, adds a second, ultralocal, auxiliary field for the same quartic, and compares both with the closed-string $\Sigma_2$.
Only the scalar model is in the paper; the Yang--Mills and stub material is only here.

\paragraph{The scalar model.}
The block below is the paper's appendix in the paper's notation.
Its statements are elementary linear algebra of a model with constant couplings and are derived in full in the block; they carry no separate entry among the results of this document.
The correspondence with the string construction is stated in the block itself, in the paragraph on the relation to the string field: the products $\phi_i \phi_j$ become $\Psi^2$, the tensor $V_{ijk\ell}$ becomes $\gs^2 \vertex{4}$, and the null vectors of $V_{IJ}$ play the role of the kernel of $\sprod{2}{2}$ (\cref{p:sec:hs:quartic}).

\subsection{Hubbard--Stratonovich transformation for \texorpdfstring{$N$}{N} scalar fields}
\label{p:app:scalar-hs}


In this appendix, we perform the Hubbard--Stratonovich transformation for a theory of several scalar fields with a quartic interaction.
We see on this example how the auxiliary fields are organized when there are several physical fields: the auxiliary field carries a pair of indices, as $\Sigma_2$ carries two strings.
This is the model of the paper's introduction with constant couplings, which underlies the analogy used in \cref{p:sec:hs:quartic}.


\subparagraph{Action of the scalar fields.}

Consider $N$ scalar fields $\phi_i$, $i = 1, \ldots, N$, with the action:
\begin{equation}
\label{p:eq:shs:action}
S
  = \int \dd^d x \left[
    - \frac{1}{2} \, \partial_\mu \phi_i \, \partial^\mu \phi_i
    - \frac{1}{2} \, M_{ij} \phi_i \phi_j
    + \frac{1}{3!} \, W_{ijk} \phi_i \phi_j \phi_k
    + \frac{1}{4!} \, V_{ijk\ell} \phi_i \phi_j \phi_k \phi_\ell
  \right],
\end{equation}
where repeated indices are summed over.
The mass matrix $M_{ij}$ and the coupling tensors $W_{ijk}$ and $V_{ijk\ell}$ are constant and completely symmetric.
We write $S_4$ for the quartic term.


\subparagraph{Auxiliary fields.}

The quartic term is the square of the products $\phi_i \phi_j$ with respect to the tensor $V_{ijk\ell}$.
We introduce one auxiliary field for each such product, that is $N (N + 1) / 2$ fields $\sigma_{ij} = \sigma_{ji}$, and add to the action the term~\cite{Stratonovich:1957:MethodCalculatingQuantum, Hubbard:1959:CalculationPartitionFunctions, Weinberg:2005:QuantumTheoryFields-2}:
\begin{equation}
\label{p:eq:shs:DeltaS4}
\Delta S_4
  := - \frac{1}{4!} \int \dd^d x \, V_{ijk\ell} \, \big(
    \sigma_{ij} - \phi_i \phi_j
  \big) \big(
    \sigma_{k\ell} - \phi_k \phi_\ell
  \big),
\end{equation}
which vanishes when $\sigma_{ij} = \phi_i \phi_j$.
We expand the square and use $V_{ijk\ell} = V_{k\ell ij}$ for the cross term, so that the quartic term of the action becomes:
\begin{equation}
\label{p:eq:shs:S4tilde}
\caction[4]
  := S_4 + \Delta S_4
  = \int \dd^d x \left[
    - \frac{1}{4!} \, V_{ijk\ell} \, \sigma_{ij} \sigma_{k\ell}
    + \frac{1}{12} \, V_{ijk\ell} \, \sigma_{ij} \phi_k \phi_\ell
  \right].
\end{equation}
The term quartic in $\phi_i$ has disappeared.
The kinetic, mass and cubic terms are untouched, so that the action
\begin{equation}
\label{p:eq:shs:action-tilde}
\begin{aligned}
  \caction
    := S + \Delta S_4
    = \int \dd^d x \bigg[
      &
      - \frac{1}{2} \, \partial_\mu \phi_i \, \partial^\mu \phi_i
      - \frac{1}{2} \, M_{ij} \phi_i \phi_j
      + \frac{1}{3!} \, W_{ijk} \phi_i \phi_j \phi_k
      \\ &
      + \frac{1}{12} \, V_{ijk\ell} \, \sigma_{ij} \phi_k \phi_\ell
      - \frac{1}{4!} \, V_{ijk\ell} \, \sigma_{ij} \sigma_{k\ell}
    \bigg]
\end{aligned}
\end{equation}
is cubic in $(\phi_i, \sigma_{ij})$.
For a single field, $N = 1$, the integrand of the quartic term~\eqref{p:eq:shs:S4tilde} reads:
\begin{equation}
- \frac{1}{4!} \, V \sigma^2
+ \frac{1}{12} \, V \sigma \phi^2.
\end{equation}
This is the transformation of the introduction at constant coupling, with the normalization $1 / 4!$ of the vertices.

The fields $\sigma_{ij}$ have no kinetic term, and their equation of motion is algebraic:
\begin{equation}
\label{p:eq:shs:eom-sigma}
V_{ijk\ell} \, \big( \sigma_{k\ell} - \phi_k \phi_\ell \big)
  = 0.
\end{equation}
When $V_{ijk\ell}$ is invertible as a matrix acting on the symmetric pairs of indices $(ij)$ and $(k\ell)$ (this is the matrix $V_{IJ}$ defined below), the solution is:
\begin{equation}
\label{p:eq:shs:solution}
\sigma_{ij}
  = \phi_i \phi_j,
\end{equation}
on which $\Delta S_4$ vanishes and $\caction$ reduces to the original action $S$.


\subparagraph{Pair-index basis.}

The pair of indices of $\sigma_{ij}$ can be traded for a single index.
Let $\lambda_{I;ij}$, $I = 1, \ldots, N (N + 1) / 2$, be an orthonormal and complete basis of the real symmetric $N \times N$ matrices:
\begin{equation}
\label{p:eq:shs:basis}
\lambda_{I;ij} \lambda_{J;ij}
  = \delta_{IJ},
\qquad
\lambda_{I;ij} \lambda_{I;k\ell}
  = \frac{1}{2} \, \big( \delta_{ik} \delta_{j\ell} + \delta_{i\ell} \delta_{jk} \big).
\end{equation}
The right-hand side of the completeness relation is the projector on symmetric matrices.
We define:
\begin{equation}
\label{p:eq:shs:pair-index}
\sigma_I
  := \lambda_{I;ij} \, \sigma_{ij},
\qquad
V_{IJ}
  := \lambda_{I;ij} \lambda_{J;k\ell} \, V_{ijk\ell},
\qquad
V_{I;k\ell}
  := \lambda_{I;ij} \, V_{ijk\ell}.
\end{equation}
Since $\sigma_{ij}$ and $V_{ijk\ell}$ are symmetric in each pair of indices, the completeness relation gives $\sigma_{ij} = \lambda_{I;ij} \sigma_I$ and $V_{I;k\ell} = V_{IJ} \lambda_{J;k\ell}$.
The quartic term~\eqref{p:eq:shs:S4tilde} then becomes:
\begin{equation}
\label{p:eq:shs:S4tilde-pair}
\caction[4]
  = \int \dd^d x \left[
    - \frac{1}{4!} \, V_{IJ} \, \sigma_I \sigma_J
    + \frac{1}{12} \, V_{I;k\ell} \, \sigma_I \phi_k \phi_\ell
  \right].
\end{equation}

The matrix $V_{IJ}$ is real and symmetric, so we can diagonalize it by an orthogonal change of the basis $\lambda_{I;ij}$, which puts the quadratic term of the $\sigma_I$ in diagonal form.
We assume that all its eigenvalues are nonzero.
Otherwise, the auxiliary fields $\sigma_I$ along the eigenvectors with zero eigenvalue appear neither in the quadratic term nor in the coupling to $\phi_i$, since $V_{I;k\ell} = V_{IJ} \lambda_{J;k\ell}$.
They do not enter the action, so we leave them out and perform the transformation on the subset of the $\sigma_I$ that does enter.


\subparagraph{Four-point amplitude.}

We read the Feynman rules off $\caction$ as in \cref{p:sec:hs:quartic}: a vertex factor is the derivative of the action with respect to the fields on its legs, and a propagator is minus the inverse of the Hessian.
The auxiliary sector~\eqref{p:eq:shs:S4tilde-pair} contributes the propagator of $\sigma_I$ and the vertex coupling $\sigma_I$ to $\phi_k \phi_\ell$, respectively:
\begin{equation}
\label{p:eq:shs:feynman-rules}
12 \, (V^{-1})_{IJ},
\qquad
\frac{1}{3!} \, V_{I;k\ell}.
\end{equation}

The exchange of $\sigma$ between two such vertices, with the external fields $\phi_i, \phi_j$ on one vertex and $\phi_k, \phi_\ell$ on the other, gives:
\begin{equation}
\label{p:eq:shs:channel}
\frac{1}{3!} \, V_{I;ij} \;
12 \, (V^{-1})_{IJ} \;
\frac{1}{3!} \, V_{J;k\ell}
  = \frac{1}{3} \, \lambda_{K;ij} V_{KI} (V^{-1})_{IJ} V_{JL} \lambda_{L;k\ell}
  = \frac{1}{3} \, V_{ijk\ell},
\end{equation}
where the first equality uses $V_{I;ij} = V_{IK} \lambda_{K;ij}$ and the second $V_{ijk\ell} = \lambda_{K;ij} V_{KL} \lambda_{L;k\ell}$, which follows from the completeness relation.
We sum over the three ways of distributing the four external fields on the two vertices and use the complete symmetry of $V_{ijk\ell}$.
The four-point amplitude is then:
\begin{equation}
\label{p:eq:shs:A4}
A_{ijk\ell}
  = \frac{1}{3} \, \big( V_{ijk\ell} + V_{ikj\ell} + V_{i\ell jk} \big) + \cdots
  = V_{ijk\ell} + \cdots,
\end{equation}
where the dots are the $s$-, $t$- and $u$-channel exchanges of $\phi$ built from the cubic coupling $W_{ijk}$, which the transformation leaves untouched.
The exchange of $\sigma$ thus reproduces the contact vertex $V_{ijk\ell}$ of the original action~\eqref{p:eq:shs:action}.


\subparagraph{Relation to the string field.}

We obtain the construction of \cref{p:sec:hs:quartic} from this one by replacing the products $\phi_i \phi_j$ by the $2$-string state $\Psi^2$, the tensor $V_{ijk\ell}$ by $\gs^2 \vertex{4}$ (so that $V_{IJ}$ becomes $\gs^2 \sprod{2}{2}$), the auxiliary field $\sigma_{ij}$ by $\Sigma_2 \in \Hn{2}$, and the propagator $12 \, (V^{-1})_{IJ}$ by $12 \gs^{-2} \, \sprod{2}{2}^{-1}$.



\subparagraph{Degenerate coupling matrix.}

To couple the directions of zero eigenvalue of $V_{IJ}$, left out above, we now add a quintic term to the action~\eqref{p:eq:shs:action}, with a constant and completely symmetric coupling $U_{ijk\ell m}$:
\begin{equation}
\label{p:eq:shs:S5}
S_5
  = \frac{1}{5!} \int \dd^d x \, U_{ijk\ell m} \, \phi_i \phi_j \phi_k \phi_\ell \phi_m.
\end{equation}
As in \cref{p:sec:hs:quintic}, we replace two pairs of fields by the auxiliary field, which turns the quintic term into a cubic coupling:
\begin{equation}
\label{p:eq:shs:S5tilde}
\caction[5]
  := \frac{1}{5!} \int \dd^d x \, U_{IJm} \, \sigma_I \sigma_J \phi_m,
\qquad
U_{IJm}
  := \lambda_{I;ij} \lambda_{J;k\ell} \, U_{ijk\ell m},
\end{equation}
which agrees with $S_5$ at $\sigma_{ij} = \phi_i \phi_j$.
When the matrix $V_{IJ}$ is invertible, eliminating $\sigma_I$ from $\caction + \caction[5]$ gives back $S + S_5$ up to terms of second order in $U_{ijk\ell m}$, the exchange of $\sigma$ between two quintic couplings, whose string counterpart is cancelled at the sextic order (\cref{p:sec:hs:sextic}).

We now suppose that $V_{IJ}$ has null vectors, $V_{IJ} \xi^J = 0$, and we work in the basis $\lambda_{I;ij}$ that diagonalizes it.
The index $I$ then splits into the directions of nonzero eigenvalue, which span the orthogonal complement $C$ of the null space, and the null directions, which span the null space $K$.
We write $\pi_C$ and $\pi_K$ for the projections on $C$ along $K$ and on $K$ along $C$.
The null components of $\sigma_I$ drop out of the quartic terms~\eqref{p:eq:shs:S4tilde-pair}, as noted above, but they enter the quintic coupling~\eqref{p:eq:shs:S5tilde}.
Under the shift $\sigma_I \to \sigma_I + \xi^I$, $\caction[5]$ changes by:
\begin{equation}
\label{p:eq:shs:null-shift}
\frac{1}{5!} \int \dd^d x \, \big(
  2 \, U_{IJm} \, \xi^I \sigma_J \phi_m
  + U_{IJm} \, \xi^I \xi^J \phi_m
\big).
\end{equation}
For this to vanish for all $\sigma_I$, the term linear in $\sigma_J$ must vanish, which then implies that the quadratic one vanishes.
The null components therefore decouple if and only if the quintic coupling annihilates the null vectors, and they can then be left out as before:
\begin{equation}
\label{p:eq:shs:kernel-condition}
\forall \xi \in K, \;
\forall k, \ell, m:
\quad
\xi^I \, U_{I;k\ell m}
  = 0,
\qquad
U_{I;k\ell m}
  := \lambda_{I;ij} \, U_{ijk\ell m}.
\end{equation}
This is the condition~\eqref{p:eq:hs5:kernel-condition} of the string construction, with $\xi$ for $N_2$ and $U_{ijk\ell m}$ for $\vertex{5}$.

When this condition fails, the null components of $\sigma_I$ have no quadratic term in $\caction$ but they have a source, the term of~\eqref{p:eq:shs:S5tilde} linear in them.
  We could keep the null components as fields, but their quadratic form in $\caction[5]$ is proportional to $\phi_m$, so that around $\phi_i = 0$ they have neither a propagator nor a perturbative expansion.
  If we leave them out instead, that is, if we take $\sigma_I$ in $C$, we get an action of the form~\eqref{p:eq:shs:action-tilde} with $V_{IJ}$ restricted to $C$.
  Since $V_{IJ}$ is invertible on $C$, $\sigma_I$ can be eliminated as before.
  The eliminated action reproduces the quartic term.
  However, at first order in $U_{ijk\ell m}$, it contains $U_{IJm} \, (\pi_C P)^I (\pi_C P)^J \phi_m$ in place of the quintic term $U_{IJm} \, P_I P_J \phi_m$, where $P_I := \lambda_{I;ij} \, \phi_i \phi_j$.
  The missing terms are those with a null direction in at least one of the two pair indices; they depend on the complement and cannot vanish for all complements.
  Indeed, we could choose any other complement of $K$ in place of the span of the eigenvectors, and we would get a different eliminated action.

We can repair the construction with a multiplier on the null directions, as in the hybrid form of \cref{p:sec:hs:sextic}.
  We keep all the components of $\sigma_I$, and we add to the action a field $\theta_I$ on the null directions, paired with the null component of $\sigma_I - P_I$:
  \begin{equation}
  \label{p:eq:shs:multiplier}
  \caction + \caction[5]
    + \int \dd^d x \, \theta_I \, \big( \pi_K (\sigma - P) \big)^I.
\end{equation}
  Since $\theta_I P_I$ is cubic, every term has at most three fields and the action stays cubic.
  The equation of motion of $\theta_I$ sets the null components of $\sigma_I$ to those of $P_I$.
  The equation of motion of the null components fixes $\theta_I$, which appears nowhere else.
  Finally, the equation of motion of the components along $C$ is the one of the action with $\sigma_I$ in $C$, with a source shifted by the quintic coupling evaluated on the null component of $P_I$.
  Eliminating all the fields gives back $S + S_5$ up to terms of second order in $U_{ijk\ell m}$, for every choice of complement.
  Indeed, the null components supply the part of $P_I$ that the complement was missing.
  The multiplier vanishes on the solution if and only if the quintic coupling annihilates the null vectors, in which case it is not needed.
  We see on this degenerate case how the multiplier replaces the quotient when a vertex does not vanish on the kernel of the quadratic form, in a setting where the linear algebra is explicit.
  For the string vertices, a multiplier on the null directions is the fallback for an accidental null vector of $\sprod{2}{2}$, though we have no indication that one exists (\cref{p:sec:hs:quintic}).


\subsection{Yang--Mills strictification and the open-string stub}
\label{sec:qft:ov}

The next block is taken from an earlier overview note.
It contains two results, with the following status.
\begin{itemize}
  \item The Yang--Mills strictification toy: integrating out the rank-three auxiliary field $G$ reproduces the non-local quartic vertex of the colour--kinematics-dual form with the correct coefficient, $\frac{1}{4} J (1 / \square) J$.
    \status{Verified}\record{001-cubic-sft-overview-note (R4)}
    The check is done on the single-contraction scalar representative $\mathscr L(G) = - \frac{1}{2} G \square G - \frac{1}{\sqrt 2} G J$, with $\square$ invertible away from zero momentum; the index structure of the rank-three auxiliary field and the strictified Lagrangian are quoted from~\cite{Borsten:2021:DoubleCopyHomotopy}, not re-derived.
  \item The open-string stub: the propagator of the auxiliary field of the stub theory, $(1 - \e^{- 2 \lambda L_0}) / L_0$ times $b_0$, is regular as $L_0 \to 0$, with limit $2 \lambda$, so that it has no on-shell pole.
    \status{Verified}\record{001-cubic-sft-overview-note (R3)}
    The check concerns the scalar factor of the propagator, with $b_0$ dropped.
    The inference from the absence of a pole to the absence of physical states in the auxiliary sector is the one of~\cite{ErbinFirat2023}, quoted, not re-derived.
    In that theory the auxiliary field lives in one copy of the open-string state space, not in a tensor product as $\Sigma_2$ does.
\end{itemize}
The remarks on the transfer to closed strings and the paragraph on the localization of the quartic potential in the superstring theory of~\cite{MaccaferriMerlano2019} are comparisons with the literature, not results, and carry no status.
The note uses the older notation of the project, translated below.

\begin{quote}
  \small
  \textbf{Translation box: overview note $\to$ paper.}
  \begin{itemize}
    \setlength{\itemsep}{0.3ex}
    \item The \enquote{HS step} is the Hubbard--Stratonovich step at the quartic order, \cref{p:sec:hs:quartic}; $\ell_3$ is the bracket of the $L_\infty$ algebra that encodes the quartic vertex $\vertex{4}$.
    \item In the stub paragraph, $\Sig \in \Hilb$ is the auxiliary field of the open-string theory with stubs, a field in the open-string state space; it is not the paper's $\Sigma_2$.
      The closing paragraph's \enquote{closed $\Sig \in \Hilb^{\otimes 2}$} is the paper's $\Sigma_2 \in \Hn{2}$; the $(\ker L_0)^c$ restriction there is the one of the stub theory, in which the auxiliary field lives outside the kernel of $L_0$.
    \item $\Mean{\cdot, \cdot}$ is the open-string BPZ product and $g_o$ the open-string coupling; $*$ is Witten's star product.
    \item $\phi_1$, $\phi_2$ are the first components of an $L_\infty$ quasi-isomorphism, and $\IBL$ is the involutive Lie bialgebra structure of \cref{ch:ibl}.
    \item \emph{Reality.} The block uses no reality condition. The note's convention elsewhere is $\Sig^\ddagger = \Sig$; on a $k$-string tensor whose strands have ghost number two, $\ddagger = (-1)^k \star$, so that $\Sig_k^\ddagger = \Sig_k$ is the paper's $\Sigma_k^\star = (-1)^k \Sigma_k$: the two agree for $\Sigma_2$, while for odd $k$ the sign is explicit, $\Sigma_3^\star = - \Sigma_3$, and $\Psi^\ddagger = \Psi$ is $\Psi^\star = - \Psi$ (\cref{ch:conventions}).
  \end{itemize}
\end{quote}

\subsubsection{The Yang--Mills strictification toy}
\label{ov:sec:explicit:YM}

The cleanest finite-dimensional illustration of the mechanism is the strictification of the
quartic Yang--Mills vertex \cite{Borsten:2021:DoubleCopyHomotopy}. Written in colour--kinematics-dual form the
quartic vertex is non-local,
\begin{equation}
	\mathscr{L}^{(4)}\sim
	-\eta^{\mu\nu}\eta^{\kappa\rho}\eta^{\lambda\sigma}\,
	\partial^{12}_\mu\partial^{34}_\nu\,
	\frac{\tr\{[A_\kappa,A_\lambda][A_\rho,A_\sigma]\}}{\square_{12}}.
	\label{ov:eq:YM4}
\end{equation}
Introducing a rank-3 auxiliary $G^a_{\mu\nu\kappa}$ (ghost number $0$, self-conjugate kinetic
term) gives the strictly cubic Lagrangian
\begin{equation}
	\mathscr{L}^{\,\mathrm{YM,st}}_4 = \tr\Big\{
	-\tfrac12\,G^{\mu\nu\kappa}_a\square G^a_{\mu\nu\kappa}
	-g\Big(\partial_\mu A_\nu+\tfrac{1}{\sqrt2}\partial^\kappa G^a_{\kappa\mu\nu}\Big)
	[A^\mu,A^\nu]\Big\}.
	\label{ov:eq:YMst}
\end{equation}
The $G$ equation of motion is $\square G^a_{\mu\nu\kappa}=-\tfrac{g}{\sqrt2}\partial_{[\mu}
[A_\nu,A_{\kappa]}]$, i.e. $G^a_{\mu\nu\kappa}=-\tfrac{g}{\sqrt2\,\square}\partial_{[\mu}
[A_\nu,A_{\kappa]}]$; substituting back reproduces the non-local quartic \eqref{ov:eq:YM4} (plus
the standard $g$-vertex). On a single-contraction scalar representative
$\mathscr{L}(G)=-\tfrac12 G\square G-\tfrac{1}{\sqrt2}GJ$ one finds the EOM
$G=-\tfrac{1}{\sqrt2}J/\square$ and the on-shell value $\mathscr{L}=\tfrac14 J(1/\square)J$, the
non-local quartic with the correct coefficient (checked symbolically). The map
$\phi_1(A)=A^{\mathrm{st}}$, $\tfrac12\phi_2(A,A)=G|_{\text{on-shell}}=
\tfrac{g}{2\square}\partial_\mu[A_\nu,A_\kappa]$ is the leading quasi-isomorphism. The auxiliary
field is purely auxiliary because $\square$ is invertible (away from zero momentum) and there
are no closed loops of purely auxiliary lines, so integrating it out is an exact Gaussian and
the theory is quantum-equivalent to ordinary Yang--Mills. This is the precise gauge-theory
mirror of the HS step: $\ell_3$ (or here $\mathscr{L}^{(4)}$) traded for cubic vertices plus an
invertible auxiliary kinetic term.

\subsubsection{The Erbin--F{\i}rat open-string stub as the \texorpdfstring{$\Aoo$}{A-infinity}
precedent}
\label{ov:sec:explicit:stub}

The closest \emph{string-theoretic} precedent is the open-string stub construction
\cite{ErbinFirat2023}, which is the $\Aoo$ analogue of the closed $\Loo$ programme. Witten's
cubic open SFT augmented with stubs of length $\lambda$ (a non-polynomial deformation) is recast
as a strictly cubic action by a single auxiliary field $\Sig\in\Hilb$ (\emph{the same} Hilbert
space as $\Psi$, not a tensor product). The gauge-invariant cubic action is
\begin{equation}
	-g_o^2 S[\Psi,\Sig]
	= \tfrac12\Mean{\Psi,Q_B\Psi}
	+ \tfrac12\Mean{\Sig,\frac{Q_B}{1-\e^{-2\lambda L_0}}\Sig}
	- \tfrac13\Mean{\Sig-\e^{-\lambda L_0}\Psi,
	(\Sig-\e^{-\lambda L_0}\Psi)*(\Sig-\e^{-\lambda L_0}\Psi)},
	\label{ov:eq:stubaction}
\end{equation}
with two independent gauge parameters $\Lambda_1,\Lambda_2\in\Hilb$. The auxiliary field is
truly auxiliary: its kinetic operator $Q_B/(1-\e^{-2\lambda L_0})$ inverts to the propagator
\begin{equation}
	\widetilde\Delta_\Sig = b_0\,\frac{1-\e^{-2\lambda L_0}}{L_0}
	= 2\lambda b_0\big[1-\lambda L_0+\cdots\big],
	\label{ov:eq:stubprop}
\end{equation}
which is \emph{regular} as $L_0\to0$ --- no on-shell pole (checked symbolically) ---
so the $\Sig$ sector has empty cohomology and carries no physical states. The key structural
facts transferring to the closed case are: (i) the construction is realised as homotopy transfer
from a \emph{strictly associative} cyclic $\Aoo$-algebra on the doubled space
$\Hilb'=\Hilb\times\Hilb$, with a \emph{strong} deformation retract
($h^2=0$, $h\iota=0$, $ph=0$) --- the ingredient not established in
\cite{SchnablStettinger2023}; (ii) integrating $\Sig$ out perturbatively
reproduces the non-polynomial stubbed theory, while the algebraic gauge choice
$\Sig=-2\sinh(\lambda L_0)\Psi$ reproduces Witten's cubic theory after a field redefinition.

The transfer to closed strings is partial. What carries over: the structural form (an auxiliary
field with an unconventional, pole-free kinetic term; homotopy transfer of a cyclic algebra on a
doubled space; the exchange of $\Sig$ carrying the region of moduli that the cubic channels leave uncovered). What is genuinely harder
\cite{ErbinFirat2023}: (a) the closed-string collar moduli are \emph{complex} (length $+$ twist),
so there is no purely real $\e^{-\lambda L_0}$ shift operator; (b) collar widths are mutually
constrained, so the binary-tree decomposition does not generalise directly; (c) $\Sig$ must live
in $\Hilb^{\otimes2}$ (or a level-matched twisted version), complicating the $(\ker L_0)^c$
restriction and the kinetic term; (d) the gauge algebra is $\Loo$, with graded-symmetric
brackets and higher gauge parameters.

\subparagraph{An additional precedent.} The EKS open-superstring localization of
Maccaferri--Merlano \cite{MaccaferriMerlano2019} provides a related explicit quartic structure:
the localized quartic effective potential factorises into matter two-point functions of
composite auxiliary fields, $\tr[\Mean{\mathbb{H}_1^{(+)}|\mathbb{H}_1^{(-)}}
+\tfrac14\Mean{\mathbb{H}_0|\mathbb{H}_0}]$, controlled by the projector condition
$P_0M_2(\Psi,\Psi)=0$.
Here the auxiliary fields live in a \emph{single} copy of $\Hilb$, not in $\Hilb^{\otimes2}$,
because the localized quartic is already a product of single-copy bilinears. This is therefore a
contrast to --- not a direct precedent for --- the closed $\Sig\in\Hilb^{\otimes2}$ structure,
and a useful warning that the tensor-product structure of the closed auxiliary field is
specific to the closed $\Loo/\IBL$ setting.



\paragraph{Two Yang--Mills mechanisms.}
The last block is the corresponding section of a note on pairings and gradings.
It repeats the colour--kinematic elimination of the previous block with the intermediate steps shown, then introduces an antisymmetric auxiliary field $h_{\mu\nu}$ with an algebraic kinetic term, whose elimination gives the Yang--Mills quartic $- \frac{1}{4} \tr([A_\mu, A_\nu] [A^\mu, A^\nu])$.
The colour--kinematic part is the toy of the previous item (\status{Verified}, same record); the ultralocal mechanism has no entry of its own among the results, and its elimination is an algebraic identity shown in full in the block.
The comparison table and the closing paragraphs are a discussion: they explain why the colour--kinematic mechanism is the better model for the invertibility of $\sprod{2}{2}$, and why the Yang--Mills analogy motivates, but does not prove, that the quartic form has no kernel on the symmetric square.

\begin{quote}
  \small
  \textbf{Translation box: pairing-grading note $\to$ paper.}
  \begin{itemize}
    \setlength{\itemsep}{0.3ex}
    \item $\mc V_4$ read as a bilinear form on $2$-string states is the paper's $\vertex{4}(A, B) = \pair[2]{A}{\sprod{2}{2}(B)}$; $p_{2,2}$ is the paper's $\sprod{2}{2}$, the same symbol.
    \item $\Sigma$ is $\Sigma_2$; $\Sym^2 \mc H_0$ is $\Hn{2}$ restricted to level-matched states, and $\Sym^2 \mc H_0 / \ker \mc V_4$ is the quotient $\Hn{2} / \ker \sprod{2}{2}$ of~\eqref{p:eq:hs4:quotient}.
    \item \enquote{The string action's $- 1 + 2 = + 1$} refers to the coefficients $- \frac{\gs^2}{4!}$ and $\frac{2 \gs^2}{4!}$ of~\eqref{p:eq:hs4:S4tilde}.
    \item \emph{Reality.} The block uses no reality condition; the note elsewhere writes $\Sigma^\ddagger = \Sigma$, which for $\Sigma_2$ is the paper's $\Sigma_2^\star = \Sigma_2$, and for odd $k$ the paper's condition carries the sign, $\Sigma_k^\star = - \Sigma_k$, with $\Psi^\star = - \Psi$ (\cref{ch:conventions}).
  \end{itemize}
\end{quote}

\subsection{Two auxiliary-field mechanisms in Yang--Mills theory}
\label{pg:sec:ym}

The finite-dimensional mirror of the Hubbard--Stratonovich step is the quartic Yang--Mills vertex, which can be made cubic with an auxiliary field in at least two ways.
The two mechanisms eliminate to the same quartic but differ in the anatomy of the auxiliary, and the comparison is what isolates the meaning of \enquote{invertible kinetic operator} for $p_{2,2}$.
We work both out, then compare.

\subsubsection{The color--kinematic mechanism}
\label{pg:sec:ym:ck}

In the color--kinematics-dual form of the Yang--Mills action the quartic vertex is written non-locally \cite{Borsten:2021:DoubleCopyHomotopy},
\begin{equation}
  \label{pg:eq:ym:quartic-nonlocal}
  \mathscr L^{(4)}
    \sim - \eta^{\mu\nu} \eta^{\kappa\rho} \eta^{\lambda\sigma} \,
      \pd^{12}_\mu \pd^{34}_\nu \,
      \frac{\tr\big( [A_\kappa, A_\lambda] [A_\rho, A_\sigma] \big)}{\square_{12}},
\end{equation}
where $\pd^{12}$ and $\square_{12}$ act on the first pair of fields.
The $1/\square$ is removed by a rank-three auxiliary field $G^a_{\mu\nu\kappa}$ with kinetic term $-\frac12 G\square G$ and a local coupling of the form $\pd A\,[A,A]$, giving a strictly cubic Lagrangian.
The mechanism is visible on a scalar representative with a single contraction, $G$ a scalar and $J$ standing for the local cubic current:
\begin{equation}
  \mathscr L(G)
    = - \frac12 \, G \square G
      - \frac{1}{\sqrt 2} \, G J .
\end{equation}
The equation of motion is $-\square G-J/\sqrt2=0$, so
\begin{equation}
  G
    = - \frac{1}{\sqrt 2} \, \frac{1}{\square} \, J,
\end{equation}
and substituting back, integrating by parts to let $\square$ act symmetrically,
\begin{equation}
  \mathscr L(G) \big|_{\rm on-shell}
    = - \frac12 \cdot \frac12 \, J \frac{1}{\square} J
      + \frac12 \, J \frac{1}{\square} J
    = \frac14 \, J \frac{1}{\square} J ,
\end{equation}
the non-local quartic with the correct coefficient.
The auxiliary kinetic operator is $\square$, invertible away from $p^2=0$.

\subsubsection{The ultralocal mechanism}
\label{pg:sec:ym:h}

The second mechanism introduces a Lie-algebra-valued antisymmetric auxiliary $h_{\mu\nu}=-h_{\nu\mu}$ with an algebraic kinetic term, no derivatives at all:
\begin{equation}
  \label{pg:eq:ym:Lh}
  \mathscr L_h
    = \frac14 \, \tr\big( h_{\mu\nu} h^{\mu\nu} \big)
      - \frac12 \, \tr\big( h_{\mu\nu} [A^\mu, A^\nu] \big).
\end{equation}
The $h$ equation of motion is algebraic,
\begin{equation}
  \frac{\pd \mathscr L_h}{\pd h_{\mu\nu}}
    = \frac12 \, h_{\mu\nu} - \frac12 \, [A_\mu, A_\nu]
    = 0
  \qquad\Longrightarrow\qquad
  \boxed{\; h_{\mu\nu} = [A_\mu, A_\nu] \;}
\end{equation}
so $h$ is exactly the commutator.
Substituting back,
\begin{equation}
  \begin{aligned}
    \mathscr L_h \big|_{h = [A, A]}
      &
      = \frac14 \, \tr\big( [A_\mu, A_\nu] [A^\mu, A^\nu] \big)
        - \frac12 \, \tr\big( [A_\mu, A_\nu] [A^\mu, A^\nu] \big)
    \\ &
      = - \frac14 \, \tr\big( [A_\mu, A_\nu] [A^\mu, A^\nu] \big),
  \end{aligned}
\end{equation}
the Yang--Mills quartic with the correct coefficient.
The kinetic operator is $\pd^2\mathscr L_h/\pd h^2=\frac12\,\mathbb 1$, a constant mass matrix, invertible everywhere with no momentum dependence.
Note the same $1:2$ magnitude ratio between the kinetic and the source term as in \eqref{p:eq:hs4:S4tilde}: at the saddle the two terms combine as $+1-2=-1$, in units of $\frac14\tr([A_\mu,A_\nu][A^\mu,A^\nu])$ --- the same doubling mechanism as the string action's $-1+2=+1$, with the overall sign flipped because the $h$-mechanism's kinetic term has the opposite sign from the $\Sigma$ one.
The elimination was checked symbolically with $\mathfrak{su}(2)$-valued $A_\mu$ in Euclidean signature; the identity is algebraic and signature-independent.

\subsubsection{The comparison}
\label{pg:sec:ym:compare}

\begin{center}
  \small
  \begin{tabular}{lccc}
    \toprule
    mechanism & auxiliary & kinetic operator & coupling \\
    \midrule
    color--kinematic & $G^a_{\mu\nu\kappa}$ (rank 3) & $\square$ (non-local) & $\pd A [A, A]$ \\
    ultralocal & $h_{\mu\nu}$ (rank 2) & $\frac12 \mathbb 1$ (ultralocal) & $h [A, A]$ \\
    closed string & $\Sigma \in \Sym^2 \mc H_0 / \ker \mc V_4$ & $\mc V_4$ (contact) & $\mc V_4(\Sigma, \Psi^2)$ \\
    \bottomrule
  \end{tabular}
\end{center}

What each mechanism buys is different.
The ultralocal one shows what \enquote{auxiliary} means at its simplest: an algebraic equation of motion, a constant invertible Hessian, no propagator, no question to ask.
The color--kinematic one shows that the kinetic operator of an auxiliary field can have a non-trivial inverse and that the substantive question is then its invertibility on the relevant space: for $G$, invertibility of $\square$ away from $p^2=0$.

The closed-string $\Sigma$ is like the $h$ form in that its kinetic operator $\mc V_4$ is of contact type, ultralocal in the moduli-space sense: the auxiliary propagates through a genuine contact vertex, with no $1/L_0^+$ and no propagator pole.
It is like the $G$ form in that invertibility is the substantive question: for $\Sigma$, invertibility of $\mc V_4$ on $\Sym^2\mc H_0$, or on the quotient by its kernel.
This is why the color--kinematic mechanism is the better model for the invertibility of $p_{2,2}$, and why it is kept rather than replaced by the simpler $h$ form: the $h$ form makes the trivially invertible end of the spectrum concrete, the $G$ form the non-trivially invertible end, and $\mc V_4$ sits between them, contact-type but with an invertibility that has to be established.

What does not carry over is the way the kernel is handled.
In Yang--Mills the operator $\square$ does have a kernel, the massless shell, and integrates out benignly through a Green's function because the degeneration is a measure-zero shell in a continuous momentum spectrum that diagonalizes $\square$.
The closed-string condition is an operator-kernel question on $\Sym^2\mc H_0$: there is no single invariant that diagonalizes the contact vertex on two-string states, and a normalizable symmetric null vector of $\mc V_4$ would be an honest flat direction of the action, not a removable shell.
The Yang--Mills analogy therefore motivates the expectation that $\mc V_4$ is non-degenerate on the symmetric square, and provides no argument for it; the quotient~\eqref{p:eq:hs4:quotient} is what makes the expectation unnecessary at this order.


